% ===============================================================
%  جلسه یازدهم
% ===============================================================
\chapter{گشتاورها، تابع مولد گشتاور، تابع مشخصه و توابع یک متغیر تصادفی}

\section{گشتاورهای یک توزیع \lr{(Moments of a Distribution)}}

\begin{definition}[گشتاور \lr{(Moment)}]
برای $n\in\mathbb W=\{0,1,2,\ldots\}$، گشتاورهای یک متغیر تصادفی مانند $X\sim f_X(x)$ به صورت زیر تعریف می‌شوند:
\[
m_n\triangleq\E(X^n)=\int_{-\infty}^{\infty}x^nf_X(x)\dd x .
\]
\textbf{گشتاورهای مرکزی} (حول میانگین $\mu$):
\[
\mu_n\triangleq\E\big((X-\mu)^n\big).
\]
همچنین گشتاورهای \textcolor{red!70!black}{مطلق} ($\E|X|^n$) و \textcolor{red!70!black}{تعمیم‌یافته} ($\E[(X-c)^n]$ حول نقطه دلخواه $c$، یا توان‌های غیرصحیح) نیز وجود دارند.
\end{definition}

\subsection{خواص}
\begin{itemize}
\item گشتاور صفرم همواره برابر با ۱، و گشتاور اول (در صورت وجود) برابر با \textbf{میانگین توزیع} هستند: $m_0=\int f_X=1$ و $m_1=\E X=\mu$. همچنین $\mu_1=0$ و $\mu_2=\Var(X)$.
\item تمام و یا برخی از گشتاورها می‌توانند وجود نداشته باشند (\textcolor{red!70!black}{توزیع کوشی})، گاهی اوقات نیز تمامی گشتاورها وجود دارند (\textcolor{blue!60!black}{گاوسی، نمایی و تمامی توزیع‌های \lr{finite-support}}).
\item تحت شرایطی، کل توزیع $f_X$ را می‌توان صرفاً با دانستن $m_0,m_1,m_2,\ldots$ به صورت یکتا به دست آورد.
\item همواره داریم:
\[
m_{2n}\ge m_n^2 .
\]
\end{itemize}
\begin{proof}[اثبات نامساوی آخر]
$\Var(X^n)=\E(X^{2n})-\big(\E X^n\big)^2\ge0$. (یا از جنسن با تابع محدب $g(u)=u^2$ روی $U=X^n$.)
\end{proof}

\begin{note}
\notetitle اگر گشتاور مرتبه $k$ وجود داشته باشد، همه گشتاورهای مرتبه کمتر نیز وجود دارند. برای توزیع‌های با تکیه‌گاه کران‌دار (مثلاً $|X|\le M$) همه گشتاورها وجود دارند زیرا $|m_n|\le M^n$. برای توزیع کوشی حتی $m_1$ وجود ندارد.
\end{note}

\subsection{گشتاور دوم، واریانس و \lr{Kurtosis}}
دقت کنید که برای یک متغیر تصادفی، سه مقدار زیر با یکدیگر فرق دارند:
\begin{itemize}
\item \textbf{گشتاور دوم:} $m_2=\E(X^2)$؛
\item \textbf{واریانس:} $\Var(X)=\E(X^2)-\E(X)^2$؛
\item \textbf{انحراف معیار:} $\sigma(X)=\sqrt{\Var(X)}$.
\end{itemize}
(این سه فقط وقتی $\E X=0$ باشد به هم مربوط و $m_2=\Var X$ می‌شود.)

\begin{definition}[چولگی و کشیدگی]
یکی از آماره‌های پرکاربرد در پردازش سیگنال (\lr{Signal Processing})، \textcolor{red!70!black}{\lr{\textbf{Kurtosis}}} (کشیدگی) نام دارد:
\[
\mathrm{Kurt}(X)=\frac{\mu_4}{\sigma^4}.
\]
به طور مشابه، \textbf{چولگی} (\lr{Skewness}) به صورت $\frac{\mu_3}{\sigma^3}$ تعریف می‌شود و میزان عدم تقارن توزیع را می‌سنجد.
\end{definition}

از \lr{Kurtosis} در \lr{Independent Component Analysis (ICA)} و جداسازی کور منابع (\lr{Blind Source Separation}) استفاده می‌گردد. برای توزیع گاوسی $\mathrm{Kurt}=3$ است؛ توزیع‌های با دم سنگین‌تر (مثل لاپلاس با $\mathrm{Kurt}=6$) کشیدگی بیشتر و توزیع‌های «پهن» (مثل یکنواخت با $\mathrm{Kurt}=1.8$) کشیدگی کمتری دارند. ایده \lr{ICA}: طبق قضیه حد مرکزی، مخلوط چند منبع مستقل «گاوسی‌تر» از تک‌تک منابع است؛ پس با بیشینه کردن فاصله \lr{Kurtosis} از ۳ می‌توان منابع را از هم جدا کرد (مثلاً جدا کردن صدای چند گوینده که هم‌زمان ضبط شده‌اند).

\begin{example}[گشتاورهای توزیع یکنواخت]
گشتاورهای $X\sim\mathcal U(0,1)$ و \lr{Kurtosis} آن را بیابید.
\begin{solution}
$m_n=\int_0^1x^n\dd x=\frac{1}{n+1}$. پس $\mu=\frac12$ و $\sigma^2=\frac13-\frac14=\frac1{12}$.
گشتاور مرکزی چهارم: $\mu_4=\int_{-1/2}^{1/2}y^4\dd y=\frac{2}{5}\cdot\frac{1}{32}=\frac1{80}$. پس
$\mathrm{Kurt}=\frac{1/80}{1/144}=\frac{144}{80}=1.8$.
\end{solution}
\end{example}

\section[تابع مولد گشتاور \lr{(MGF)}]{تابع مولد گشتاور {\large\lr{(Moment Generating Function)}}}

\begin{definition}[تابع مولد گشتاور]
تابع مولد گشتاور (\lr{MGF}) برای متغیر تصادفی $X\sim f_X(x)$ و $t\in\mathbb C$ به صورت زیر تعریف می‌شود:
\[
M_X(t)\triangleq\E\big[e^{tX}\big]=\int_{-\infty}^{\infty}e^{tx}f_X(x)\dd x .
\]
محدوده تعریف ورودی تابع \lr{MGF}، تمامی مقادیری از $t\in\mathbb C$ است که مقدار میانگین فوق به ازای آن وجود داشته باشد.
\end{definition}

\subsection{خواص}
\begin{itemize}
\item تابع مولد گشتاور به ازای هر $t\in\R$ (صرفاً مقادیر حقیقی) همواره مقداری مثبت است (چون $e^{tX}>0$).
\item تابع \lr{MGF} در اثر میانگین‌گیری بر روی یک ماهیت تصادفی (یعنی $e^{tX}$) به وجود آمده است. لذا \lr{MGF} تصادفی نیست و تابعی یقینی به حساب می‌آید.
\item تمام اطلاعات توزیع $f_X(x)$ در \lr{MGF} آن وجود دارد. به عبارت دیگر، تابع \lr{MGF} حاصل انتقال توزیع $f_X(x)$ به \textcolor{red!70!black}{یک حوزه دیگر (مشابه حوزه لاپلاس)} است. (در واقع $M_X(t)$ همان تبدیل لاپلاس دوطرفه $f_X$ در $s=-t$ است.)
\item می‌توان به راحتی دید که $M_X(0)=1$. همچنین برای \textcolor{blue!60!black}{توزیع گاوسی}، \lr{MGF} به ازای تمامی صفحه اعداد مختلط وجود دارد.
\end{itemize}

\subsection{کاربرد اصلی: تولید گشتاورها}
می‌توان به راحتی دید که خاصیت زیر برای مشتقات \lr{MGF} وجود دارد (به دلیل خطی بودن عملگر مشتق، از انتگرال عبور خواهد کرد):
\[
\frac{\partial^n}{\partial t^n}M_X(t)=\frac{\partial^n}{\partial t^n}\int_{-\infty}^{\infty}e^{tx}f_X(x)\dd x=\int_{-\infty}^{\infty}x^ne^{tx}f_X(x)\dd x .
\]
در صورتی که مشتق $n$ام تابع \lr{MGF} را با $M_X^{(n)}(t)$ نشان دهیم، خواهیم داشت:
\[
\boxed{\ M_X^{(n)}(0)=\int_{-\infty}^{\infty}x^nf_X(x)\dd x=m_n\ }
\]
یعنی مشتق $n$ام در صفر، همان \textcolor{blue!60!black}{گشتاور} مرتبه $n$ برای متغیر تصادفی $X$ است. معادلاً با بسط تیلور:
\[
M_X(t)=\E\Big[\sum_{n\ge0}\frac{t^nX^n}{n!}\Big]=\sum_{n\ge0}\frac{m_n}{n!}t^n ;
\]
\lr{MGF} همه گشتاورها را به صورت ضرایب سری تیلورش در خود «بسته‌بندی» کرده است ــ از این‌رو «مولد گشتاور» نام دارد.

\begin{example}[\lr{MGF} و گشتاورهای توزیع نمایی]
\begin{solution}
برای $t<\lambda$:
$M_X(t)=\int_0^\infty e^{tx}\lambda e^{-\lambda x}\dd x=\lambda\int_0^\infty e^{-(\lambda-t)x}\dd x=\frac{\lambda}{\lambda-t}$ (برای $t\ge\lambda$ انتگرال واگراست).
مشتقات: $M'(t)=\frac{\lambda}{(\lambda-t)^2}$ و $M''(t)=\frac{2\lambda}{(\lambda-t)^3}$؛ پس $m_1=\frac1\lambda$ و $m_2=\frac{2}{\lambda^2}$ و $\Var X=\frac{1}{\lambda^2}$ (مطابق جلسه قبل). در حالت کلی $M_X(t)=\frac{1}{1-t/\lambda}=\sum_n\frac{t^n}{\lambda^n}$، پس $m_n=\frac{n!}{\lambda^n}$.
\end{solution}
\end{example}

\begin{example}[\lr{MGF} توزیع گاوسی]
\begin{solution}
ابتدا برای $Z\sim\mathcal N(0,1)$ با «تکمیل مربع» $tz-\frac{z^2}{2}=-\frac{(z-t)^2}{2}+\frac{t^2}{2}$:
\[
M_Z(t)=\int\frac{1}{\sqrt{2\pi}}e^{tz-z^2/2}\dd z=e^{t^2/2}\int\frac{1}{\sqrt{2\pi}}e^{-(z-t)^2/2}\dd z=e^{t^2/2}.
\]
برای $X=\mu+\sigma Z$: $M_X(t)=\E e^{t\mu+t\sigma Z}=e^{t\mu}M_Z(\sigma t)=e^{t\mu+\sigma^2t^2/2}$.
\end{solution}
\end{example}

\subsection{خاصیت مهم: \lr{MGF} جمع متغیرهای مستقل}
\begin{theorem}
در صورتی که $X,Y$ دو متغیر تصادفی مستقل با توابع مولد گشتاور $M_X(t),M_Y(t)$ باشند، برای متغیر تصادفی جدید $X+Y$ خواهیم داشت:
\[
\begin{aligned}
M_{X+Y}(t)&=\E\big[e^{t(X+Y)}\big]=\E\big[e^{tX+tY}\big]=\E\big[e^{tX}e^{tY}\big]\\
&=\E\big[e^{tX}\big]\E\big[e^{tY}\big]=M_X(t)\,M_Y(t).
\end{aligned}
\]
(در سطر دوم از استقلال $e^{tX}$ و $e^{tY}$ و قاعده $\E(UV)=\E U\,\E V$ استفاده شد.)
\end{theorem}
این خاصیت بسیار قدرتمند است: توزیع جمع متغیرهای مستقل (که مستقیماً با «کانولوشن» به دست می‌آید) در حوزه \lr{MGF} به یک ضرب ساده تبدیل می‌شود ــ درست مثل تبدیل لاپلاس/فوریه در سیگنال‌ها و سیستم‌ها.

\begin{example}[جمع پواسون‌های مستقل]
اگر $X\sim\mathrm{Pois}(\lambda_1)$ و $Y\sim\mathrm{Pois}(\lambda_2)$ مستقل باشند، توزیع $X+Y$ چیست؟
\begin{solution}
$M_X(t)=\sum_ke^{tk}e^{-\lambda}\frac{\lambda^k}{k!}=e^{-\lambda}e^{\lambda e^t}=e^{\lambda(e^t-1)}$. پس
$M_{X+Y}(t)=e^{\lambda_1(e^t-1)}e^{\lambda_2(e^t-1)}=e^{(\lambda_1+\lambda_2)(e^t-1)}$ که \lr{MGF} توزیع $\mathrm{Pois}(\lambda_1+\lambda_2)$ است. طبق یکتایی (پایین‌تر)، $X+Y\sim\mathrm{Pois}(\lambda_1+\lambda_2)$. مثلاً اگر به دو صف مستقل به ترتیب ۲ و ۳ مشتری در ساعت برسد، به مجموع آن‌ها پواسون با نرخ ۵ می‌رسد. به همین روش، جمع دو گاوسی مستقل نیز گاوسی است: $e^{t\mu_1+\sigma_1^2t^2/2}e^{t\mu_2+\sigma_2^2t^2/2}$ = \lr{MGF} $\mathcal N(\mu_1+\mu_2,\sigma_1^2+\sigma_2^2)$.
\end{solution}
\end{example}

\subsection{\lr{MGF} توزیع‌های مشهور}
\begin{center}
\renewcommand{\arraystretch}{1.6}
\begin{tabular}{l|l}
\textbf{توزیع} & \textbf{تابع مولد گشتاور} $M_X(t)$\\\hline
تباهیده (ثابت) $\delta_a$ & $e^{ta}$\\
برنولی $P(X=1)=p$ & $1-p+pe^t$\\
دوجمله‌ای $B(n,p)$ & $(1-p+pe^t)^n$\\
هندسی $(1-p)^kp$ ($k\ge0$) & $\dfrac{p}{1-(1-p)e^t},\quad t<-\ln(1-p)$\\
دوجمله‌ای منفی $\mathrm{NB}(r,p)$ & $\Big(\dfrac{p}{1-e^t+pe^t}\Big)^r,\quad t<-\ln(1-p)$\\
پواسون $\mathrm{Pois}(\lambda)$ & $e^{\lambda(e^t-1)}$\\
یکنواخت پیوسته $U(a,b)$ & $\dfrac{e^{tb}-e^{ta}}{t(b-a)}$\\
یکنواخت گسسته $DU(a,b)$ & $\dfrac{e^{at}-e^{(b+1)t}}{(b-a+1)(1-e^t)}$\\
لاپلاس $L(\mu,b)$ & $\dfrac{e^{t\mu}}{1-b^2t^2},\quad |t|<1/b$\\
نرمال $N(\mu,\sigma^2)$ & $e^{t\mu+\sigma^2t^2/2}$\\
نمایی $\mathrm{Exp}(\lambda)$ & $\dfrac{\lambda}{\lambda-t},\quad t<\lambda$
\end{tabular}
\end{center}

\begin{note}
\notetitle دوجمله‌ای = جمع $n$ برنولی مستقل، پس \lr{MGF} آن توان $n$ام \lr{MGF} برنولی است. به همین ترتیب دوجمله‌ای منفی (تعداد شکست‌ها تا $r$امین موفقیت) جمع $r$ هندسی مستقل است و \lr{MGF} آن توان $r$ام \lr{MGF} هندسی است.
\end{note}

\section{تابع مشخصه \lr{(Characteristic Function)}}

\begin{definition}[تابع مشخصه]
تابع مشخصه یا \lr{Characteristic Function} برای $\omega\in\R$ به صورت زیر تعریف می‌شود:
\[
\Phi(\omega)=M_X(j\omega)=\int_{-\infty}^{\infty}e^{j\omega x}f_X(x)\dd x .
\]
\end{definition}
\begin{itemize}
\item همان‌گونه که مشخص است، این تابع به نوعی بیانگر \textbf{تبدیل فوریه} تابع چگالی احتمال متغیر تصادفی است (با علامت مثبت در نما).
\item عمده خواص \lr{MGF}، از جمله تولید گشتاور، برای \lr{Characteristic Function} هم وجود دارد: $\Phi^{(n)}(0)=j^nm_n$ و برای متغیرهای مستقل $\Phi_{X+Y}=\Phi_X\Phi_Y$.
\item مزیت بزرگ: چون $|e^{j\omega x}|=1$، انتگرال همواره همگراست؛ پس تابع مشخصه \textbf{برای هر توزیعی} (حتی کوشی) وجود دارد، در حالی که \lr{MGF} ممکن است وجود نداشته باشد.
\end{itemize}

\subsection{خواص و قضیه یکتایی}
\begin{theorem}
\begin{enumerate}
\item تابع مشخصه در $\omega=0$ بیشترین مقدار قدرمطلق (\lr{Magnitude}) را داراست:
\[
|\Phi(\omega)|\le\Phi(0)=1 .
\]
\item \textbf{(یکتایی)} در صورتی که برای دو متغیر تصادفی $X,Y$ داشته باشیم $\Phi_X(\omega)=\Phi_Y(\omega)$ (برای همه $\omega$)، آنگاه دو متغیر تصادفی هم‌توزیع (\lr{identically distributed}) هستند.
\end{enumerate}
\end{theorem}
\begin{proof}[اثبات بند ۱]
$|\Phi(\omega)|=\big|\int e^{j\omega x}f_X(x)\dd x\big|\le\int|e^{j\omega x}|f_X(x)\dd x=\int f_X(x)\dd x=1=\Phi(0)$.
\end{proof}
(بند ۲ نتیجه وارون‌پذیری تبدیل فوریه است: $f_X(x)=\frac{1}{2\pi}\int\Phi(\omega)e^{-j\omega x}\dd\omega$.)

\subsection{توابع مشخصه توزیع‌های معروف}
\begin{center}
\renewcommand{\arraystretch}{1.8}
\begin{tabular}{l|l}
گاوسی $\mathcal N(\mu,\sigma^2)$ & $\Phi_X(\omega)=e^{j\mu\omega-\sigma^2\omega^2/2}$\\
نمایی $\mathrm{Exponential}(\lambda)$ & $\Phi_X(\omega)=\Big(1-\dfrac{j\omega}{\lambda}\Big)^{-1}$\\
یکنواخت $\mathrm{Uniform}(a,b)$ & $\Phi_X(\omega)=e^{j\omega\frac{a+b}{2}}\,\dfrac{\sin\big((b-a)\omega/2\big)}{(b-a)\omega/2}$\\
برنولی $\mathrm{Bern}(p)$ & $\Phi_X(\omega)=pe^{j\omega}+(1-p)$\\
کوشی $(x_0,\gamma)$ & $\Phi_X(\omega)=e^{jx_0\omega-\gamma|\omega|}$
\end{tabular}
\end{center}
\begin{note}
\notetitle (در اسلاید، فرمول یکنواخت بدون ضریب $\frac12$ داخل سینوس آمده؛ شکل درست همان است که در جدول بالا نوشته‌ایم: $\int_a^b\frac{e^{j\omega x}}{b-a}\dd x=\frac{e^{j\omega b}-e^{j\omega a}}{j\omega(b-a)}=e^{j\omega\frac{a+b}2}\frac{2\sin(\omega\frac{b-a}2)}{\omega(b-a)}$.)

تابع مشخصه \textcolor{red!70!black}{توزیع گاوسی}، خود به شکل تابع چگالی احتمال یک \textcolor{red!70!black}{توزیع گاوسی} است ($e^{-\sigma^2\omega^2/2}$ یک «زنگوله» است) ــ همان خاصیت معروف «تبدیل فوریه گاوسی، گاوسی است». همچنین کوشی در نقطه $\omega=0$ مشتق‌پذیر نیست (به دلیل $|\omega|$)، که با عدم وجود میانگین آن سازگار است.
\end{note}

\section{توابعی از یک متغیر تصادفی}

\subsection{تعریف مسئله}
فرض کنید که متغیر تصادفی $X\sim f_X(x)$ و تابع دلخواه $g:\R\to\R$ داده شده باشند. متغیر تصادفی جدیدی به شکل $Y=g(X)$ را تشکیل می‌دهیم. هدف، به دست آوردن \textcolor{red!70!black}{تابع چگالی احتمال $f_Y(y)$} بر حسب $f_X(x)$ و $g(\cdot)$ است.

\begin{center}
\begin{tikzpicture}[>=Stealth]
\draw[thick] plot[smooth cycle,tension=0.8] coordinates {(0,0) (1.2,0.9) (2.6,0.6) (3,-0.4) (2,-1.3) (0.5,-1.1)};
\fill[gray!40] (1.5,-0.1) ellipse (0.75 and 0.45);
\node at (1.5,-0.1) {$X\in E$};
\node at (0.4,-0.8) {$\Omega$};
\draw[->,thick] (6,-1.4) -- (6,1.4) node[above]{}; \node at (6.3,-1.3) {$\R$};
\fill[gray!60] (5.93,-0.3) rectangle (6.07,0.5); \node at (6.3,0.1) {$E$};
\draw[->,thick] (9.5,-1.4) -- (9.5,1.4); \node at (9.8,-1.3) {$\R$};
\fill[gray!60] (9.43,0.6) rectangle (9.57,1.0); \node[anchor=west,font=\scriptsize] at (9.65,0.8) {$\{g(x):x\in E\}$};
\draw[->,thick] (2.2,0.2) to[out=50,in=150] node[above]{$X$} (5.9,0.2);
\draw[->,thick] (6.1,0.3) to[out=50,in=150] node[above]{$g$} (9.4,0.8);
\draw[->,thick] (2.1,-0.4) to[out=-40,in=-130] node[below]{$Y=g(X)$} (9.4,0.6);
\end{tikzpicture}
\end{center}

\textbf{روش کلی (از طریق \lr{CDF}):}
\[
F_Y(y)=\Prob(g(X)\le y)=\Prob\big(X\in\{x:\ g(x)\le y\}\big)=\int_{\{x:\,g(x)\le y\}}f_X(x)\dd x,
\]
و سپس $f_Y=F_Y'$. این روش برای هر تابع $g$ کار می‌کند.

\subsection{موارد ساده}
\textbf{شیفت.} فرض کنید که متغیر تصادفی $Y$ به صورت $Y=X+b$ تعریف شده باشد، که در اینجا $b$ یک مقدار ثابت و یقینی است. در این صورت، هر مقدار تصادفی $Y$ از طریق ابتدا نمونه‌گیری از $X$ با استفاده از $f_X(x)$، و سپس شیفت دادن آن به اندازه مقدار ثابت و یقینی $b$ حاصل می‌شود. پس می‌توان به راحتی استدلال کرد که:
\[
\boxed{\ f_Y(y)=f_X(y-b)\ }
\]
(از طریق \lr{CDF}: $F_Y(y)=\Prob(X\le y-b)=F_X(y-b)$ و با مشتق‌گیری حکم به دست می‌آید.)

\textbf{تبدیل خطی.} حال فرض کنید که متغیر تصادفی $Y$ به صورت $Y=aX+b$ تعریف شده باشد، که در اینجا $a\ne0,b$ مقادیر ثابت و یقینی هستند. می‌توان نوشت (برای $a>0$):
\[
\begin{aligned}
f_Y(y)&=\lim_{h\to0}\frac{\Prob(y\le aX+b\le y+h)}{h}=\lim_{h\to0}\frac{\Prob(y-b\le aX\le y+h-b)}{h}\\
&=\lim_{h\to0}\frac{\Prob\big(\frac{y-b}{a}\le X\le\frac{y-b}{a}+\frac{h}{|a|}\big)}{h}
=\frac{1}{|a|}\lim_{h\to0}\frac{\Prob\big(\frac{y-b}{a}\le X\le\frac{y-b}{a}+\frac{h}{|a|}\big)}{h/|a|}\\
&=\frac{1}{|a|}f_X\Big(\frac{y-b}{a}\Big).
\end{aligned}
\]
(برای $a<0$ جهت نامساوی‌ها برعکس می‌شود و همان نتیجه با $|a|$ به دست می‌آید.)
\[
\boxed{\ f_Y(y)=\frac{1}{|a|}f_X\Big(\frac{y-b}{a}\Big)\ }
\]

\begin{note}
\notetitle \textbf{تعبیر:} ضریب $\frac{1}{|a|}$ از آنجاست که تبدیل، بازه‌ها را $|a|$ برابر کش می‌دهد؛ برای حفظ مساحت (احتمال) کل، ارتفاع چگالی باید $|a|$ برابر کم شود. به عنوان نتیجه مهم: اگر $X\sim\mathcal N(\mu,\sigma^2)$ آنگاه $aX+b\sim\mathcal N(a\mu+b,a^2\sigma^2)$؛ به‌ویژه $Z=\frac{X-\mu}{\sigma}\sim\mathcal N(0,1)$ (استانداردسازی).
\end{note}

\begin{example}[تبدیل خطی یکنواخت]
اگر $X\sim\mathcal U(0,1)$ و $Y=3-2X$، چگالی $Y$ را بیابید.
\begin{solution}
$a=-2$، $b=3$: $f_Y(y)=\frac12f_X\big(\frac{y-3}{-2}\big)=\frac12$ برای $0\le\frac{3-y}{2}\le1$، یعنی $1\le y\le3$. پس $Y\sim\mathcal U(1,3)$.
\end{solution}
\end{example}

\begin{example}[یک تبدیل غیرخطی: $Y=X^2$]
اگر $X\sim\mathcal N(0,1)$، چگالی $Y=X^2$ را بیابید.
\begin{solution}
برای $y>0$: $F_Y(y)=\Prob(-\sqrt y\le X\le\sqrt y)=F_X(\sqrt y)-F_X(-\sqrt y)$. با مشتق‌گیری (قاعده زنجیره‌ای):
\[
f_Y(y)=\frac{1}{2\sqrt y}\big[f_X(\sqrt y)+f_X(-\sqrt y)\big]=\frac{1}{\sqrt{2\pi y}}e^{-y/2},\qquad y>0 .
\]
این توزیع «کای‌دو با یک درجه آزادی» ($\chi^2_1$) نام دارد. دقت کنید که چون $g(x)=x^2$ یک‌به‌یک نیست، دو شاخه $\pm\sqrt y$ هر دو سهم دارند. (این ایده در جلسه بعد به فرمول کلی تبدیل تعمیم داده می‌شود.)
\end{solution}
\end{example}

\section{تمرین‌های تکمیلی}

\begin{exercise}
با استفاده از \lr{MGF}، میانگین و واریانس توزیع دوجمله‌ای را به دست آورید.
\end{exercise}
\begin{solution}
$M(t)=(1-p+pe^t)^n$. $M'(t)=npe^t(1-p+pe^t)^{n-1}$، پس $M'(0)=np$.
$M''(t)=npe^t(\cdots)^{n-1}+n(n-1)p^2e^{2t}(\cdots)^{n-2}$، پس $M''(0)=np+n(n-1)p^2$. واریانس: $np+n(n-1)p^2-n^2p^2=np(1-p)$.
\end{solution}

\begin{exercise}
اگر $M_X(t)=\frac{1}{4}+\frac12e^{t}+\frac14e^{2t}$، توزیع $X$ را مشخص کنید.
\end{exercise}
\begin{solution}
$M_X(t)=\sum_kP(X=k)e^{tk}$، پس با مقایسه ضرایب: $P(X=0)=\frac14$، $P(X=1)=\frac12$، $P(X=2)=\frac14$ ــ یعنی $\mathrm{Bin}(2,\frac12)$ (تعداد شیر در دو پرتاب). در واقع $M_X(t)=\big(\frac12+\frac12e^t\big)^2$.
\end{solution}

\begin{exercise}
اگر $X\sim\mathrm{Exp}(1)$، چگالی $Y=2X+1$ و $\E Y$ را بیابید.
\end{exercise}
\begin{solution}
$f_Y(y)=\frac12e^{-(y-1)/2}$ برای $y\ge1$ (نمایی با نرخ $\frac12$ که یک واحد شیفت خورده است). $\E Y=2\E X+1=3$.
\end{solution}

\begin{exercise}
نشان دهید \lr{Kurtosis} توزیع گاوسی برابر ۳ است.
\end{exercise}
\begin{solution}
برای $Z\sim\mathcal N(0,1)$: $M_Z(t)=e^{t^2/2}=1+\frac{t^2}{2}+\frac{t^4}{8}+\cdots$. ضریب $t^4$ برابر $\frac{m_4}{4!}$ است، پس $m_4=\frac{24}{8}=3$. چون کشیدگی به تبدیل خطی حساس نیست ($\mu_4$ و $\sigma^4$ هر دو در $\sigma^4$ ضرب می‌شوند)، برای هر گاوسی $\mathrm{Kurt}=3$.
\end{solution}
